\documentclass[12pt,a4paper]{article}
\usepackage{geometry}
\geometry{a4paper, margin=1in}
\usepackage{setspace}
\usepackage{hyperref}
\usepackage{xepersian}

\settextfont[Scale=1.2]{XB Niloofar}

\linespread{1.6}

\begin{document}

\begin{titlepage}
\begin{center}
\vspace*{1.5cm}
{\huge\bfseries دانشکده علوم ریاضی} \\
\vspace{0.3cm}
{\Large دانشگاه فردوسی مشهد} \\
\vspace{2.5cm}
{\Huge\bfseries پروژه درس پایگاه داده} \\
\vspace{0.5cm}
{\Large طراحی پایگاه داده سامانه مسیریابی شهری} \\
\vspace{3.5cm}
{\Large استاد درس: دکتر جلال‌الدین نصیری} \\
\vspace{1.5cm}
{\large تیم حل تمرین:} \\
\vspace{0.3cm}
{\large محمد باغبان‌زاده، یگانه راد} \\
\vspace{3.5cm}
{\large بهار ۱۴۰۵}
\end{center}
\end{titlepage}

\newpage
\tableofcontents
\newpage

\section{مقدمه و توصیف پروژه}
فرض کنید قرار است پایگاه داده یک سامانه مسیریابی شهری طراحی شود. این سامانه قرار است نقشه شهر را نمایش دهد و برای حرکت درون شهر، مسیرهای مناسب را پیشنهاد کند. انتخاب مسیر باید با در نظر گرفتن شکل شهر، مسیرهای موجود، قوانین عبور و مرور، شرایط زمانی و نوع وسیله نقلیه انجام شود. تمرکز این پروژه بر طراحی پایگاه داده‌ای است که بتواند این اطلاعات را به شکل درست و قابل تحلیل ذخیره کند، نه پیاده‌سازی الگوریتم‌های مسیریابی.

ابتدای یا انتهای یک مسیر، شهر به مجموعه‌ای از نقاط میانی تقسیم می‌شود. هر نقطه می‌تواند یک تقاطع خیابان یا نقطه‌ای قابل استفاده در فرآیند مسیریابی باشد. هر نقطه دارای موقعیت جغرافیایی مشخص است. برخی نقاط ممکن است در برخی بازه‌های زمانی قابل استفاده نباشند، برای مثال به دلیل مسدود شدن مسیر یا شرایط خاص. این وضعیت باید در داده‌ها قابل ثبت باشد تا در زمان پیشنهاد مسیر در نظر گرفته شود.

بین این نقاط، مسیرهایی وجود دارد که امکان حرکت را فراهم می‌کنند. هر مسیر یک بخش پیوسته از شبکه معابر شهر است و دو نقطه را به هم متصل می‌کند. مسیرها ممکن است یک‌طرفه یا دوطرفه باشند. هر مسیر دارای ویژگی‌هایی است که در انتخاب مسیر اهمیت دارند، مانند طول مسیر یا نوع کلی آن.

حرکت در مسیرها همیشه آزاد نیست و قوانین مختلفی می‌توانند استفاده از آن‌ها را محدود کنند. برخی از این قوانین به زمان وابسته هستند و فقط در ساعات یا روزهای خاصی اعمال می‌شوند. بعضی محدودیت‌ها به نوع وسیله نقلیه مربوط می‌شوند. همچنین ممکن است جهت حرکت یک مسیر در بازه‌های زمانی خاص تغییر کند. این قوانین می‌توانند در طول زمان تغییر کنند و باید به گونه‌ای ذخیره شوند که بتوان اعتبار زمانی آن‌ها را بررسی کرد.

سامانه از چند نوع وسیله نقلیه پشتیبانی می‌کند. هر وسیله نقلیه ویژگی‌های کلی خود را دارد و نمی‌تواند از همه مسیرها استفاده کند. حتی اگر استفاده از یک مسیر مجاز باشد، نحوه عبور وسایل مختلف از آن مسیر ممکن است متفاوت باشد. به همین دلیل، باید بتوان رابطه بین نوع وسیله نقلیه و مسیرها را در پایگاه داده مدل‌سازی کرد.

وضعیت حرکت در شهر در طول زمان تغییر می‌کند. یک مسیر ممکن است در ساعتی از روز خلوت و در ساعتی دیگر پرترافیک باشد. سامانه باید بتواند وضعیت‌های زمانی مسیرها را ذخیره کند و علاوه بر آن، الگوهای معمول رفتاری مسیرها در زمان‌های مختلف را نیز نگهداری کند تا امکان تحلیل‌های آینده وجود داشته باشد.

علاوه بر تغییرات معمول، رویدادهای غیرعادی نیز می‌توانند باعث اختلال در حرکت شوند. این رویدادها ممکن است باعث افزایش زمان عبور یا غیرقابل استفاده شدن یک مسیر شوند. هر رویداد دارای زمان شروع، زمان پایان تقریبی و میزان تاثیرگذاری است و ممکن است پس از بررسی تایید یا رد شود. به عنوان مثال بروز یک تصادف در یک مسیر.

در برخی شرایط، عوامل محیطی مانند وضعیت آب و هوا نیز می‌توانند بر مسیرها تاثیر بگذارند. این تاثیر برای همه مسیرها یکسان نیست و برخی مسیرها نسبت به شرایط خاص حساس‌تر هستند. بنابراین لازم است تاثیر شرایط محیطی به صورت وابسته به مسیر و زمان قابل ثبت باشد.

سامانه باید امکان استفاده از حمل و نقل عمومی را نیز در نظر بگیرد. خطوط حمل و نقل عمومی دارای مسیر مشخص و مجموعه‌ای از ایستگاه‌ها هستند که ترتیب آن‌ها اهمیت دارد. کاربران ممکن است بخواهند مسیرهایی را مشاهده کنند که ترکیبی از حرکت با وسیله شخصی و حمل و نقل عمومی باشد، بنابراین این ارتباط باید در داده‌ها قابل نمایش باشد.

\newpage
\section{ساختار پروژه}
پروژه به چندین فاز تقسیم شده است. هر فاز بر اساس مراحل طراحی و پیاده‌سازی پایگاه داده تعریف شده و باید به ترتیب انجام شود. تمرکز بر طراحی مفهومی، نرمال‌سازی، کوئری‌نویسی و جنبه‌های پیشرفته است. در ادامه، هر فاز به همراه توضیحات دقیق و انتظارات آمده است.

\subsection{فاز اول: طراحی مدل \lr{Entity-Relationship}}
در این فاز، مدل مفهومی پایگاه داده را با استفاده از دیاگرام \lr{ER} طراحی کنید. دیاگرام باید شامل تمام موجودیت‌ها (\lr{Entities})، روابط (\lr{Relationships})، ویژگی‌ها (\lr{Attributes}) و محدودیت‌ها (\lr{Constraints}) بر اساس توصیف پروژه باشد.
\begin{itemize}
\item موجودیت‌های اصلی را شناسایی کنید (مانند نقاط میانی، مسیرها، وسایل نقلیه، قوانین محدودیت، رویدادها، حمل و نقل عمومی، درخواست‌های کاربران و غیره).
\item روابط بین موجودیت‌ها را مدل کنید (مانند اتصال مسیرها به نقاط، محدودیت‌های زمانی بر مسیرها، رابطه وسایل نقلیه با مسیرها).
\item از نمادهای استاندارد \lr{ER} استفاده کنید (مانند مستطیل برای موجودیت‌ها، لوزی برای روابط، بیضی برای ویژگی‌ها).
\item محدودیت‌های کاردینالیتی (\lr{Cardinality}) و مشارکت (\lr{Participation}) را مشخص کنید.
\item اگر لازم باشد، موجودیت‌های ضعیف (\lr{Weak Entities}) یا روابط چندبه‌چند (\lr{Many-to-Many}) را با جدول‌های واسط مدل کنید.
\item دیاگرام را به صورت خوانا و با ابزارهایی مانند \lr{Draw.io} یا \lr{Lucidchart} رسم کنید.
\end{itemize}

\subsection{فاز دوم: نرمال‌سازی تا سطح مناسب}
در این فاز، مدل \lr{ER} را به جداول رابطه‌ای تبدیل کنید و سپس نرمال‌سازی را تا سطحی که برای این پروژه مناسب است انجام دهید.
\begin{itemize}
\item ابتدا مدل \lr{ER} را به \lr{schema} رابطه‌ای تبدیل کنید، شامل کلیدهای اصلی (\lr{Primary Keys})، کلیدهای خارجی (\lr{Foreign Keys}) و ویژگی‌ها.
\item وابستگی‌های تابعی (\lr{Functional Dependencies}) را شناسایی کنید.
\item نرمال‌سازی را گام‌به‌گام انجام دهید.
\item توضیح دهید چرا سطح نرمال‌سازی انتخابی مناسب است (مثلاً تعادل بین جلوگیری از ناهنجاری‌ها و کارایی کوئری‌ها).
\item اگر \lr{denormalization} در بخشی لازم باشد (مانند برای بهبود عملکرد در داده‌های زمانی)، آن را توجیه کنید.
\item \lr{schema} نهایی را به صورت لیست جداول با ستون‌ها، انواع داده‌ها (\lr{Data Types}) و محدودیت‌ها (مانند \lr{NOT NULL}، \lr{UNIQUE}، \lr{NULL}) ارائه دهید.
\end{itemize}

\subsection{فاز سوم: پاسخ به کوئری‌ها با جبر رابطه‌ای و \lr{SQL}}
در این فاز، کوئری‌های مشخص شده را با استفاده از جبر رابطه‌ای و \lr{SQL} پیاده‌سازی کنید. فرض کنید از \lr{schema} فاز دوم استفاده می‌شود و نام جداول و ستون‌ها را بر اساس آن تنظیم کنید. کوئری‌ها به دو بخش تقسیم شده‌اند:

\subsubsection{کوئری‌های جبر رابطه‌ای}
\begin{enumerate}
\item لیست تمام مسیرهایی که وسایل نقلیه از نوع "دوچرخه" می‌توانند از آن‌ها استفاده کنند را به همراه طول هر مسیر نمایش دهید.
\item تمام تقاطع‌هایی را پیدا کنید که حداقل یک محدودیت گردش (\lr{turn restriction}) دارند.
\item مسیرهایی که برای "خودرو شخصی" قابل استفاده هستند اما برای "کامیون" قابل استفاده نیستند را پیدا کنید.
\item جفت مسیرهایی را پیدا کنید که در یک تقاطع مشترک به هم متصل هستند و امکان پیچ از مسیر اول به مسیر دوم وجود ندارد (\lr{turn restriction}).
\end{enumerate}

\subsubsection{کوئری‌های \lr{SQL}}
\begin{enumerate}
\item تعداد کل ایستگاه‌های حمل و نقل عمومی در هر خط را به همراه نام خط محاسبه کنید.
\item میانگین طول مسیرهایی که وسایل نقلیه سنگین (کامیون) می‌توانند از آن‌ها عبور کنند را محاسبه کنید.
\item نقاطی که تعداد مسیرهای متصل به آن‌ها بیشتر از میانگین است را پیدا کنید. (استفاده از \lr{Window Function} الزامی است)
\item برای هر مسیر، میانگین زمان سفر در ساعات اوج ترافیک (۷--۹ صبح و ۱۷--۱۹ عصر) را محاسبه کنید و با میانگین کل مقایسه کنید.
\item همه نقاطی که از یک نقطه مشخص (مثلا \lr{node\_id = 1}) با حداکثر ۳ مسیر قابل دسترسی هستند را پیدا کنید. (استفاده از \lr{Recursive CTE} الزامی است)
\end{enumerate}

\subsection{فاز چهارم (امتیازی)}
\begin{enumerate}
\item[سوال ۱:] سامانه نیاز دارد مسیرهایی را نمایش دهد که در زمان جاری سیستم برای خودرو شخصی مجاز هستند. یک \lr{View} طراحی کنید که فقط مسیرهای مجاز برای خودرو شخصی را در زمان جاری نمایش دهد. فرض کنید قوانین عبور و مرور دارای بازه‌های زمانی اعتبار (\lr{valid\_from}، \lr{valid\_to}) هستند. توضیح دهید \lr{View} شما چگونه این بازه‌های زمانی را مدیریت می‌کند.
\item[سوال ۲:] برای هر نوع وسیله نقلیه مسیرهایی که بیشترین سرعت متوسط را دارند را همراه با رتبه‌بندی نمایش دهید. (استفاده از \lr{RANK()} یا \lr{ROW\_NUMBER()} الزامی است.)
\end{enumerate}

\section{نکات تکمیلی}
\begin{itemize}
\item پروژه به صورت گروه‌های دو نفره صورت می‌گیرد.
\item هرگونه کپی، تشابه و تقلب موجب نمره صفر می‌شود.
\item تسلط شما حین ارائه پروژه بر روی قسمت‌های پروژه و سوالاتی که از شما پرسیده می‌شود بخشی از نمره پروژه شما را تشکیل می‌دهد.
\item پروژه باید به صورت کامل و منظم مستندسازی شود، شامل توضیحات کد، دیاگرام‌ها و گزارش نهایی.
\end{itemize}

\end{document}
